% ECONOMICS_PROBLEM: {"id": "O34-376", "title": "Protection of AI outputs", "jel_code": "O34", "source_id": "OP-0416"}
% !TeX program = xelatex
\documentclass[11pt,a4paper]{article}
\usepackage[margin=23mm]{geometry}
\usepackage{fontspec}
\setmainfont{Latin Modern Roman}
\usepackage{amsmath,amssymb,mathtools}
\usepackage{xcolor,enumitem,booktabs,longtable,array,xurl,hyperref}
\definecolor{muted}{HTML}{53636C}
\hypersetup{colorlinks=true,urlcolor=blue,linkcolor=blue}
\setlength{\parindent}{0pt}
\setlength{\parskip}{5pt}
\newcommand{\R}{\mathbb R}
\newcommand{\E}{\mathbb E}
\newcommand{\N}{\mathbb N}
\newcommand{\ind}{\mathbf 1}
\newcommand{\Var}{\operatorname{Var}}
\newcommand{\Cov}{\operatorname{Cov}}
\newcommand{\supp}{\operatorname{supp}}
\DeclareMathOperator*{\argmin}{arg\,min}
\DeclareMathOperator*{\argmax}{arg\,max}
\newcommand{\entrymeta}[3]{{\sffamily\small\color{muted}\textbf{\texttt{#1}}\quad|\quad #2\par #3\par}}
\newcommand{\Problemref}[1]{\texttt{#1}}
\newcommand{\JELref}[2]{\texttt{#2} (JEL \texttt{#1})}
\begin{document}
\section*{Protection of AI outputs}
\textbf{Problem ID:} \texttt{O34-376}\quad \textbf{JEL:} \texttt{O34}\par
% BEGIN ECONOMICS STATEMENT
\label{problem:OP-0416}
{\sffamily\small\color{muted}Original subject group: Innovation and AI economics\par}
\label{definitions:problem:30007698}
\textbf{Audit classification:} Background; proposed extension. The source supplies economic IP trade-offs but does not establish this particular dynamic optimization as open. Legal-policy feasibility must not be silently asserted.
\subsection*{Definitions and precise research target}
Consider a finite-horizon market for creative works, with periods \(t=0,\ldots,H\), discount factor \(\delta\in(0,1]\), human creators and firms using AI to produce works. A protection rule \(\ell\in\mathcal L\) specifies duration, permitted reuse and enforceable licensing rights for AI-assisted outputs; \(\mathcal L\) is a finite set of legally feasible policy alternatives fixed for the exercise. Creators and firms choose investments, output prices and licensing offers in a stated dynamic equilibrium. Let \(B_t(\ell)\) be aggregate consumption value, \(C_t(\ell)\) real creation and distribution cost, and \(K_t(\ell)\) enforcement cost. Define
\[
W(\ell)=E\!\left[\sum_{t=0}^{H}\delta^t\{B_t(\ell)-C_t(\ell)-K_t(\ell)\}\right].
\]
Licensing payments are transfers within this welfare calculation; prices influence access and investment through equilibrium choices. New creators may enter, and later works may use earlier works, so cumulative innovation must be represented rather than assumed independent across periods. Characterize or estimate which \(\ell\) maximizes \(W(\ell)\) for a specified creative sector, allowing uncertainty about demand, substitution and reuse technology. Distinguish increases in output volume from changes in originality and access. This is a scoped proposed policy-design problem supported by economic IP trade-offs, rather than a canonical unsolved theorem or a claim about what protection current law provides.

\textbf{Definitions, assumptions and mathematical qualifications.} A policy rule identifies output categories, duration, permitted reuse, infringement verification and licensing constraints. The legally feasible alternatives are inputs to the model rather than claims about current law. Give a finite dynamic action/state model, shock law, entry opportunities and equilibrium concept; select an equilibrium in each regime with a stated kernel. Aggregate consumption value is gross willingness to pay at realized access, costs are resource expenditures, and enforcement transfers cancel only within the welfare boundary. \(H\in\mathbb N\), \(\delta\in(0,1]\) and finite expected payoff components make the displayed welfare finite. The finite nonempty policy menu guarantees a maximum once welfare is known.
\subsection*{Variants and assumption changes}
\begin{enumerate}
\item \textbf{Editorial, one period.} Set \(H=0\) and hold existing investments fixed to study access and substitution across protected and unprotected outputs.
\item \textbf{Editorial, cumulative innovation.} Let later works use earlier ones and allow endogenous entry/investment; identify how licensing restrictions shift both access and new work quality over the fixed horizon.
\end{enumerate}
\subsection*{References and source audit}
\textbf{Checked source.} NBER Working Paper 32698 (July 2024), indexed abstract; published chapter DOI 10.1086/732853 (2025). Primary metadata and abstract indexed; full official chapter returned HTTP 403. \url{https://www.nber.org/papers/w32698}.
\begin{enumerate}\raggedright
\item\hypertarget{definitions:source:30007698:1}{} Gaétan de Rassenfosse; Adam B. Jaffe; Joel Waldfogel. 2024. \emph{Intellectual Property and Creative Machines}. NBER Working Paper 32698 (July 2024), indexed abstract; published chapter DOI 10.1086/732853 (2025).
\url{https://www.nber.org/papers/w32698}.
DOI: \url{https://doi.org/10.3386/w32698}.
\end{enumerate}

\subsection*{Common conventions: Source mathematical conventions}

These conventions apply to each entry unless explicitly replaced.

\textbf{Models and identification.} A primitive model \(m\) specifies all probability laws, feasible choices, information, equilibrium and measurement rules used by its target. The admissible class \(\mathcal M\) is fixed independently of outcomes. If \(P_O(m)\) is its observable probability law and \(\tau(m)\) the target, its identified set at observable law \(P\) is
\[
 \mathcal I_\tau(P)=\{\tau(m):m\in\mathcal M,\ P_O(m)=P\}.
\]
Point identification means that this set is a singleton. An empty set signals incompatibility of data and assumptions. Coordinate bounds need not describe the joint set. Bounds are sharp only if every retained target is attainable or arbitrarily approachable by compatible models. Finite bounds require explicit restrictions on unobserved quantities; unrestricted confounding, hidden wealth, extrapolation or arbitrary equilibrium selection can yield uninformative sets.

\textbf{Causal interventions.} A potential outcome \(Y_i(p)\) is the outcome under a complete policy regime \(p\), including timing, financing, information and market spillovers. The target population and units are fixed before treatment. With interference, use the complete assignment vector or a justified exposure mapping; individual no-interference notation alone cannot identify an economy-wide intervention. A difference of mean potential outcomes does not require the cross-world joint distribution. Conditioning on post-treatment migration, survival, enrollment or employment changes the estimand and needs separate justification.

\textbf{Statistical statements.} An estimator, confidence set, forecast or test specifies a sampling law, model class and loss/coverage criterion. Uniform \((1-\alpha)\)-coverage means \(\inf_{m\in\mathcal M}\Pr_m\{\tau(m)\in C_n\}\geq1-\alpha\), or an explicitly asymptotic version. The whole parameter space trivially has coverage, so informativeness requires a stated width, risk or power objective. A forecasting improvement is relative to a named benchmark, information set and evaluation population.

\textbf{Equilibrium and optimization.} An equilibrium comprises feasible plans, optimality, beliefs where needed, budget constraints and all market/resource clearing. The primitives or a stated benchmark define each correspondence \(\mathcal E(p;m)\). Nonemptiness is assumed or proved, never inferred just from compact policy sets. A selected-equilibrium target uses a declared selection rule; a robust target quantifies over all permitted equilibria. Use suprema or infima unless attainment is established. An \(\varepsilon\)-optimal policy is feasible and within \(\varepsilon\) of the finite optimum value. Report an empty feasible set separately.

\textbf{Welfare and accounting.} Normative utility, interpersonal normalization, social weights, numeraire, discounting and the use of public receipts are primitives. Behavioral utility need not coincide with the welfare criterion. Household welfare includes ownership income and rents once; adding profits again to compensating variation generally double-counts them. A monetary equivalent is defined by an explicit transfer and price convention, with monotonicity and range conditions. Average effects, mechanism contributions, transfers and welfare losses are different targets. Infinite sums require convergence and dynamic feasible plans require terminal or no-Ponzi conditions.

\textbf{Source versions.} A working-paper date, revision date and journal publication year are distinguished. A DOI for a journal version is not automatically the DOI of an earlier manuscript. Locators refer to the stated version and printed pages where specified. A metadata-only check does not verify a claim about a full-text conclusion. Every entry's source subsection narrows the provenance claim accordingly.
% END ECONOMICS STATEMENT
\end{document}
